\documentclass[tikz,border=8pt]{standalone}
\usepackage{ctex}
\usepackage{amsmath}
\usetikzlibrary{calc, arrows.meta, positioning}
\begin{document}
\begin{tikzpicture}[scale=1.25, >={Stealth[length=2.4mm]}]
  % ===== 目标函数等高线（以无约束极小 U 为圆心的同心圆）=====
  \coordinate (U) at (4.16,1.53);   % 无约束极小（不可行）
  \coordinate (xs) at (2.65,1.0);   % 约束最优解 x*（边界切点）
  \foreach \r in {0.6,1.1,2.2} {
    \draw[gray!55, thin] (U) circle (\r);
  }
  % 与可行域相切的那条等高线（最低可达值）——高亮
  \draw[orange!80!black, thick] (U) circle (1.6);

  % ===== 可行域 F（凸多边形，浅蓝填充）=====
  \fill[blue!12, draw=blue!55, thick]
    (0,0) -- (3,0) -- (2.3,2) -- (0,2.6) -- cycle;
  \node[font=\small, blue!55!black] at (1.05,1.55) {可行域 $\mathcal{F}$};
  \node[font=\scriptsize, blue!55!black, align=center] at (1.15,1.12)
    {$\{x: g_i(x)\le0,\ h_j(x)=0\}$};

  % ===== 关键点 =====
  % 无约束极小 U
  \fill[gray!70!black] (U) circle (1.6pt);
  \node[font=\scriptsize, align=center, anchor=west] at ($(U)+(0.18,0.05)$)
    {无约束极小\\$\nabla f=0$（不可行）};
  % 约束最优解 x*
  \fill[red] (xs) circle (1.8pt);
  \node[font=\scriptsize, red, anchor=north east] at ($(xs)+(-0.12,-0.10)$) {约束最优解 $x^*$};

  % ===== KKT 预览：x* 处 -∇f 与约束法向同向，指向 U =====
  \draw[->, red, thick] (xs) -- ($(xs)!0.42!(U)$)
    node[midway, above, sloped, font=\scriptsize] {$-\nabla f \parallel \nabla g$};

  % ===== 等高线标注 =====
  \node[font=\scriptsize, orange!80!black, anchor=south] at (2.75,2.95)
    {目标函数 $f$ 的等高线};
  \draw[orange!80!black, thin] (2.95,2.9) -- (3.2,2.74);

  % ===== 标题 =====
  \node[font=\bfseries] at (2.4,3.5)
    {最优化问题的几何：在可行域 $\mathcal{F}$ 内最小化 $f$};

  % ===== 半透明注释框（右下空白处）=====
  \node[font=\scriptsize, fill=yellow!18, draw=yellow!60!black, rounded corners=2pt,
        align=left, fill opacity=0.9, text opacity=1, anchor=north west]
    at (3.05,0.55)
    {等高线越往里 $f$ 越小；\\能进入 $\mathcal{F}$ 的最小等高线\\恰与边界相切于 $x^*$};

\end{tikzpicture}
\end{document}
